\section*{Methods}

\subsection*{The model}
The primary experiment uses the external feedback architecture shown in Figure 1A of Sussillo and Abbott \cite{sussillo_generating_2009}. The model consists of a recurrent generator network, a linear readout unit, and a feedback pathway from the readout output back to the recurrent network. The recurrent network contains \(N\) rate units. The internal state is denoted by \(x(t)\), and the firing-rate vector is
\[
r(t)=\tanh(x(t)).
\]

The recurrent connectivity matrix \(M\) is randomly generated with connection probability \(p\) and scaling \(g/\sqrt{pN}\). The network output is a linear readout of the firing rates:
\[
z(t)=w^T r(t),
\]
where \(w\) is the readout weight vector. During the Figure 1A experiment, only \(w\) is modified; the recurrent matrix \(M\) and the feedback weights \(w_f\) remain fixed.

The discrete-time dynamics used in the Matlab implementation are
\[
x(t+\Delta t)
=
(1-\Delta t)x(t)
+
\Delta t M r(t)
+
\Delta t w_f z(t).
\]
The output \(z(t)\) is fed back to the network, allowing the trained readout to influence future network state. This closed-loop structure enables autonomous generation after training.

\subsection*{Target function}
The target is the Figure 2D signal from the paper and supplemental code:
\[
f(t)=\frac{1}{1.5}\left[
a\sin(\pi q t)+\frac{a}{2}\sin(2\pi q t)
+\frac{a}{6}\sin(3\pi q t)+\frac{a}{3}\sin(4\pi q t)
\right],
\]
with \(a=1.3\) and \(q=1/60\) for the external-feedback experiment.

\subsection*{FORCE/RLS training}
At each learning step, the instantaneous error is
\[
e(t)=z(t)-f(t).
\]
The readout vector is updated with recursive least squares:
\[
k=P r,\quad c=\frac{1}{1+r^T k},\quad
P \leftarrow P-k k^T c,\quad
w \leftarrow w-e k c.
\]
The matrix \(P\) is initialized as \(I/\alpha\). In the code this update is isolated in \texttt{code/rls\_update.m}; the full Figure 1A simulation is implemented in \texttt{code/run\_force\_external.m}.

Training and testing are separated. During training, \(w\) is updated every \texttt{learn\_every} time steps. During testing, learning is stopped and the network continues from the final training state. The output is then compared with the target over the next time interval.

\subsection*{Algorithm summary}
\begin{enumerate}
    \item Initialize random network \(M\), feedback \(w_f\), readout \(w=0\), and \(P=I/\alpha\).
    \item For each training time step, update \(x\), \(r=\tanh(x)\), and \(z=w^Tr\).
    \item At each learning step, compute \(e=z-f(t)\), then update \(w\) and \(P\) with RLS.
    \item Store the training output \(z(t)\) and the readout norm \(\|w\|_2\).
    \item Stop learning and continue the network forward to record autonomous testing output.
\end{enumerate}
